\section{La modalidad: Tesis}
Presentas y defiendes tu tesis ante tu sínodo. Es la modalidad que la plataforma toma como referencia: los oficios hablan de «su tesis».

\subsection{Conocer la modalidad y comenzar}
En \emph{Posgrado: obtener mi grado} está la lista de modalidades; cada una abre su página con los requisitos en resumen y los manuales, sin iniciar nada. Para comenzar, acepta el aviso de privacidad y pulsa \boton{Comenzar mi trámite de grado}. La modalidad se elige después, en tus datos.
\captura{tesis_egresado_01_entrada}{«Posgrado: obtener mi grado»: los cuatro momentos, las modalidades y sus manuales.}
\captura{tesis_egresado_02_informativa}{Página de la modalidad «Tesis», antes de iniciar: en qué consiste, requisitos en resumen y manuales.}
\captura{tesis_egresado_04_comenzar}{Aceptas el aviso de privacidad y pulsas «Comenzar mi trámite de grado».}
\section{Lo que cambia en «Tesis»}
El recorrido es el mismo en todas las modalidades. En los oficios y formatos el trabajo se nombra «tesis», y en el formato de integración de expediente se marca la casilla de esta modalidad.
